\documentclass{article}
\usepackage[T1]{fontenc}
\usepackage{lmodern}
\usepackage{graphicx}
\usepackage{amsmath,amssymb}
\usepackage[a4paper,margin=2cm]{geometry}
\usepackage[absolute,overlay]{textpos}
\setlength{\TPHorizModule}{1mm}
\setlength{\TPVertModule}{1mm}
\usepackage[indonesian]{babel}
\usepackage{hyperref}
\setlength{\emergencystretch}{2em}
% unit: Halaman 25

\begin{document}

% === Halaman 25 (python/native) ===

• P(B) merupakan luaran dari fungsi f. 

• Bit-bit P(B) di-XOR-kan dengan Li – 1 menghasilkan Ri:

           Ri = Li – 1  P(B)

• Jadi, luaran dari putaran ke-i adalah


(Li, Ri) = (Ri – 1 , Li – 1  P(B)) 

f



Li-1

Ri

32 bit

32 bit

25

\end{document}
